% SVG Studio — 架构信息图
% 叙事：十二个透明画布研究，将矢量几何作为可编程媒介。
% 构图：左右分层。左栏 = 项目总览（项目身份 + 核心理念 + 12 场景家族图）。
%       右栏 = 技术栈与构建流水线（渲染管道 + 测试/验证 + 目录结构）。
% 读法：从左上标题进入，沿左栏向下理解"是什么 / 有哪些场景"，
%       再沿右栏向下理解"怎么实现 / 怎么验证"。

\documentclass[border=12pt,varwidth=19cm]{standalone}
\usepackage[UTF8,fontset=none]{ctex}
\IfFontExistsTF{Source Han Serif SC}{
  \setmainfont[BoldFont={Source Han Serif SC Bold}]{Source Han Serif SC}
  \setsansfont[BoldFont={Source Han Serif SC Bold}]{Source Han Serif SC}
  \setCJKmainfont[BoldFont={Source Han Serif SC Bold}]{Source Han Serif SC}
  \setCJKsansfont[BoldFont={Source Han Serif SC Bold}]{Source Han Serif SC}
}{
  \PackageWarningNoLine{tikz-infographic}{Source Han Serif SC not found; using Fandol}
  \setCJKmainfont[BoldFont={FandolSong-Bold}]{FandolSong-Regular}
  \setCJKsansfont[BoldFont={FandolHei-Bold}]{FandolHei-Regular}
}
\IfFontExistsTF{Sarasa Mono SC}{
  \setmonofont{Sarasa Mono SC}
  \setCJKmonofont{Sarasa Mono SC}
}{
  \PackageWarningNoLine{tikz-infographic}{Sarasa Mono SC not found; using TeX defaults}
  \setmonofont{Latin Modern Mono}
  \setCJKmonofont{FandolFang-Regular}
}
\usepackage{xcolor}
\usepackage{tikz}
\usepackage{listings}
\usetikzlibrary{arrows.meta,backgrounds,calc,fit,positioning,shapes.geometric}
\pgfdeclarelayer{edge layer}
\pgfdeclarelayer{back layer}
\pgfsetlayers{background,back layer,edge layer,main}

% ── 语义调色板 ────────────────────────────────────────────────────────────
% 深色太空背景 + 霓虹色点缀，呼应项目本身的深空蓝美学。
\definecolor{Deep}{HTML}{0D1B30}       % 主背景
\definecolor{Panel}{HTML}{132741}      % 面板背景
\definecolor{Ink}{HTML}{EDF8FF}        % 主文字
\definecolor{Muted}{HTML}{82A4BA}      % 次级文字
\definecolor{Fog}{HTML}{57768C}        % 弱线条
\definecolor{Teal}{HTML}{54D6C6}       % 青绿（系统/边）
\definecolor{Blue}{HTML}{7B9CFF}       % 蓝（流/队列）
\definecolor{Coral}{HTML}{FF6F91}      % 珊瑚（产品/告警）
\definecolor{Gold}{HTML}{FFD166}       % 金（重点/标注）
\definecolor{Purple}{HTML}{B889FF}     % 紫（输出）
\definecolor{Paper}{HTML}{F7F4EE}      % 对比浅色区

% 代码样式
\lstdefinestyle{deep-code}{
  basicstyle=\ttfamily\scriptsize\color{Ink},
  keywordstyle=\bfseries\color{Gold},
  commentstyle=\itshape\color{Teal!80},
  stringstyle=\color{Coral},
  numberstyle=\scriptsize\color{Muted},
  numbers=left,
  numbersep=7pt,
  frame=single,
  rulecolor=\color{Fog},
  backgroundcolor=\color{Panel},
  breaklines=true,
  columns=fullflexible,
  keepspaces=true,
  showstringspaces=false,
  tabsize=2,
  xleftmargin=8pt,
  framexleftmargin=6pt,
  aboveskip=5pt,
  belowskip=0pt
}

\tikzset{
  every node/.style={font=\rmfamily},
  % 主色带：渲染主路径
  spine/.style={draw=Teal!35, line width=5mm, line cap=round, line join=round},
  % 家族色块
  family/.style={
    rounded corners=2.5mm, line width=1.2pt,
    minimum width=22mm, minimum height=14mm,
    align=center, font=\rmfamily\footnotesize\bfseries, text=Ink,
    inner sep=2pt
  },
  % 技术模块
  module/.style={
    rounded corners=2mm, line width=1pt,
    minimum width=30mm, minimum height=10mm,
    align=center, font=\rmfamily\footnotesize\bfseries, text=Ink,
    inner sep=3pt
  },
  % 阶段编号圆
  stage/.style={
    circle, draw=Teal, fill=Teal, text=Deep, inner sep=0pt,
    minimum size=7.5mm, font=\rmfamily\scriptsize\bfseries
  },
  % 连线
  call/.style={
    -{Latex[length=2mm,width=1.4mm]}, draw=Teal,
    line width=1pt, rounded corners=2mm
  },
  reply/.style={
    -{Latex[length=2mm,width=1.4mm]}, draw=Blue,
    densely dashed, line width=.9pt, rounded corners=2mm
  },
  flow/.style={
    -{Latex[length=2mm,width=1.4mm]}, draw=Gold,
    line width=1pt, rounded corners=2mm
  },
  boundary/.style={draw=Fog!60, dash pattern=on 2pt off 2.4pt, line width=.6pt},
  % 文字
  kicker/.style={font=\rmfamily\scriptsize\bfseries, text=Teal},
  h1/.style={font=\rmfamily\huge\bfseries, text=Ink},
  h2/.style={font=\rmfamily\large\bfseries, text=Ink},
  h3/.style={font=\rmfamily\footnotesize\bfseries, text=Gold},
  body/.style={font=\rmfamily\scriptsize, text=Muted},
  label/.style={font=\rmfamily\scriptsize\bfseries, text=Ink},
  note/.style={font=\rmfamily\scriptsize, text=Muted},
  pill/.style={
    rounded corners=3mm, fill=Panel, text=Teal,
    inner xsep=2.8mm, inner ysep=1mm,
    font=\rmfamily\scriptsize\bfseries, draw=Teal!50, line width=.6pt
  },
  % 场景小徽章
  dot/.style={circle, inner sep=0pt, minimum size=3mm},
}

\begin{document}
\setlength{\parindent}{0pt}
\begin{tikzpicture}[x=1cm,y=1cm]

  % ── 背景 ─────────────────────────────────────────────────────────────────
  \begin{pgfonlayer}{background}
    \fill[Deep,rounded corners=3mm] (-.6,-.5) rectangle (18.4,25.0);
    % 细微网格点
    \foreach \x in {0,1,...,18}
      \foreach \y in {0,1,...,24}
        \fill[Fog!18] (\x,\y) circle (.3pt);
  \end{pgfonlayer}

  % ── 标题区 ───────────────────────────────────────────────────────────────
  \node[kicker,anchor=west] at (0,24.2) {SVG STUDIO / 矢量几何工作室};
  \node[h1,anchor=west] at (0,23.4) {十二个透明画布，};
  \node[h1,anchor=west] at (0,22.5) {将矢量几何作为可编程媒介};
  \node[body,anchor=west] at (0,21.7)
    {Paper.js 构建精确几何 · Rough.js 注入手绘质感 · Playwright 做端到端像素级验证};

  % 右上：元数据徽章
  \node[pill,anchor=east] at (17.9,23.7) {12 个场景};
  \node[pill,anchor=east] at (17.9,22.8) {9 大家族};
  \node[pill,anchor=east] at (17.9,21.9) {全透明画布};

  % ── 分割线 ───────────────────────────────────────────────────────────────
  \draw[boundary] (0,21.0) -- (17.9,21.0);

  % ════════════════════════════════════════════════════════════════════════
  % 左栏：场景家族图谱（主角：12 个场景按家族组织）
  % ════════════════════════════════════════════════════════════════════════

  \node[h2,anchor=west] at (0,20.2) {场景家族图谱};
  \node[body,anchor=west] at (0,19.65)
    {每个场景回答一个设计问题；从有机几何到信息设计，从科学图示到编排图形。};

  % 三行四列的家族网格（12 个场景）
  % 行 1: botanical, metro, orbits, topology
  % 行 2: isometric-city, wave-lab, contour-map, circuit
  % 行 3: timeline, molecule, loom, type-system

  \begin{scope}[local bounding box=scenes]
    % 第一行
    \node[family, fill=Coral!18, draw=Coral!70] (s1) at (1.8,18.3) {植物花房\\[-1mm]\scriptsize\color{Muted}botanical};
    \node[family, fill=Blue!18, draw=Blue!70] (s2) at (4.8,18.3) {地铁线路\\[-1mm]\scriptsize\color{Muted}metro};
    \node[family, fill=Gold!18, draw=Gold!70] (s3) at (7.8,18.3) {星轨图\\[-1mm]\scriptsize\color{Muted}orbits};
    \node[family, fill=Teal!18, draw=Teal!70] (s4) at (10.8,18.3) {系统拓扑\\[-1mm]\scriptsize\color{Muted}topology};

    % 第二行
    \node[family, fill=Purple!18, draw=Purple!70] (s5) at (1.8,16.5) {等距城市\\[-1mm]\scriptsize\color{Muted}isometric-city};
    \node[family, fill=Coral!18, draw=Coral!70] (s6) at (4.8,16.5) {波动实验室\\[-1mm]\scriptsize\color{Muted}wave-lab};
    \node[family, fill=Teal!18, draw=Teal!70] (s7) at (7.8,16.5) {等高线地图\\[-1mm]\scriptsize\color{Muted}contour-map};
    \node[family, fill=Gold!18, draw=Gold!70] (s8) at (10.8,16.5) {模拟电路\\[-1mm]\scriptsize\color{Muted}circuit};

    % 第三行
    \node[family, fill=Blue!18, draw=Blue!70] (s9) at (1.8,14.7) {任务时间线\\[-1mm]\scriptsize\color{Muted}timeline};
    \node[family, fill=Coral!18, draw=Coral!70] (s10) at (4.8,14.7) {分子几何\\[-1mm]\scriptsize\color{Muted}molecule};
    \node[family, fill=Gold!18, draw=Gold!70] (s11) at (7.8,14.7) {算法织机\\[-1mm]\scriptsize\color{Muted}loom};
    \node[family, fill=Purple!18, draw=Purple!70] (s12) at (10.8,14.7) {字体系统\\[-1mm]\scriptsize\color{Muted}type-system};
  \end{scope}

  % 家族标注（右侧 9 大家族）
  \node[h3,anchor=west] at (13.0,19.1) {九大家族};
  \node[body,anchor=west] at (13.0,18.55) {有机几何 · 信息设计};
  \node[body,anchor=west] at (13.0,18.05) {科学图示 · 系统图};
  \node[body,anchor=west] at (13.0,17.55) {空间图 · 科学仪表盘};
  \node[body,anchor=west] at (13.0,17.05) {制图 · 技术原理图};
  \node[body,anchor=west] at (13.0,16.55) {叙事时间线 · 生成矢量};
  \node[body,anchor=west] at (13.0,16.05) {编辑图形};

  % 复杂度标签
  \node[pill,anchor=west] at (13.0,15.2) {全部：advanced 级别};

  % catalog.json 数据来源
  \node[note,anchor=west] at (0,13.9)
    {数据来源：\texttt{catalog.json} · 每条记录包含 id、用途、家族、设计问题、复杂度与标签};

  % ════════════════════════════════════════════════════════════════════════
  % 中下：渲染管道（主叙事）
  % ════════════════════════════════════════════════════════════════════════

  \node[h2,anchor=west] at (0,12.9) {渲染管道};
  \node[body,anchor=west] at (0,12.35)
    {从场景选择到透明画布输出：Paper.js 绘制精确路径，Rough.js 叠加手绘纹理。};

  % 渲染流程阶段
  \node[stage] (p1) at (1.2,11.2) {01};
  \node[stage] (p2) at (4.5,11.2) {02};
  \node[stage] (p3) at (7.8,11.2) {03};
  \node[stage] (p4) at (11.1,11.2) {04};
  \node[stage] (p5) at (14.4,11.2) {05};

  \node[label,anchor=south] at (1.2,11.75) {URL 参数\\选择场景};
  \node[label,anchor=south] at (4.5,11.75) {Paper.js\\构建几何};
  \node[label,anchor=south] at (7.8,11.75) {Rough.js\\注入手绘感};
  \node[label,anchor=south] at (11.1,11.75) {Canvas\\透明渲染};
  \node[label,anchor=south] at (14.4,11.75) {omitBackground\\导出 PNG};

  \node[note,anchor=north] at (1.2,10.65) {\texttt{?scene=botanical}};
  \node[note,anchor=north] at (4.5,10.65) {贝塞尔/路径/渐变};
  \node[note,anchor=north] at (7.8,10.65) {粗䊁度 · 弯曲};
  \node[note,anchor=north] at (11.1,10.65) {1400×900};
  \node[note,anchor=north] at (14.4,10.65) {Alpha 通道保留};

  \begin{pgfonlayer}{edge layer}
    \draw[spine] (1.2,11.2) -- (4.5,11.2);
    \draw[spine] (4.5,11.2) -- (7.8,11.2);
    \draw[spine] (7.8,11.2) -- (11.1,11.2);
    \draw[spine] (11.1,11.2) -- (14.4,11.2);
  \end{pgfonlayer}

  % 渲染就绪标志
  \node[pill,anchor=center] at (17.0,11.2) {\texttt{\_\_VIS\_READY\_\_}};
  \draw[flow] (14.4,11.2) -- (16.1,11.2);

  % ════════════════════════════════════════════════════════════════════════
  % 左下：代码片段（核心结构）
  % ════════════════════════════════════════════════════════════════════════

  \node[h3,anchor=west] at (0,9.2) {核心代码结构};
  \node[body,anchor=west] at (0,8.75)
    {\texttt{src/main.ts} 统一定义 12 个渲染函数，通过查找表分发。};

  \node[anchor=north west] at (0,8.5) {
    \begin{minipage}{10.8cm}
    \begin{lstlisting}[style=deep-code,language=c++]
const W = 1400, H = 900;
const scene = new URLSearchParams(location.search)
  .get('scene') ?? 'botanical';

function botanical() { /* 花瓣、渐变、标注 */ }
function metro()     { /* 线路、站点、Rough边框 */ }
// ... 共 12 个渲染函数

const renderers = { botanical, metro, orbits,
  topology, 'isometric-city': isometricCity,
  'wave-lab': waveLab, 'contour-map': contourMap,
  circuit, timeline, molecule, loom,
  'type-system': typeSystem };
(renderers[scene] ?? botanical)();
window.__VIS_READY__ = true;
    \end{lstlisting}
    \end{minipage}
  };

  % ════════════════════════════════════════════════════════════════════════
  % 右下：测试与验证流水线
  % ════════════════════════════════════════════════════════════════════════

  \node[h3,anchor=west] at (12.0,9.2) {像素级验证流水线};
  \node[body,anchor=west] at (12.0,8.75)
    {Playwright 无头浏览器 + pngjs 像素统计。};

  \node[module, fill=Panel, draw=Teal!60, anchor=north, align=left, text width=34mm]
    (t1) at (14.0,8.0) {
      \hphantom{0}1. Vite 构建\\
      \scriptsize\color{Muted}\hphantom{000}启动本地服务器
    };
  \node[module, fill=Panel, draw=Blue!60, anchor=north, align=left, text width=34mm]
    (t2) at (14.0,7.05) {
      \hphantom{0}2. Playwright 加载\\
      \scriptsize\color{Muted}\hphantom{000}断网模式 + headless
    };
  \node[module, fill=Panel, draw=Gold!60, anchor=north, align=left, text width=34mm]
    (t3) at (14.0,6.1) {
      \hphantom{0}3. 交互触发验证\\
      \scriptsize\color{Muted}\hphantom{000}pointer 事件计数
    };
  \node[module, fill=Panel, draw=Coral!60, anchor=north, align=left, text width=34mm]
    (t4) at (14.0,5.15) {
      \hphantom{0}4. 像素统计断言\\
      \scriptsize\color{Muted}\hphantom{000}透明 / 可见 / 彩色像素
    };

  \begin{pgfonlayer}{edge layer}
    \draw[flow] (t1.south) -- (t2.north);
    \draw[flow] (t2.south) -- (t3.north);
    \draw[flow] (t3.south) -- (t4.north);
  \end{pgfonlayer}

  % 断网标识
  \node[pill, anchor=west] at (12.0,4.4) {禁外网：仅 localhost + data: 通过};

  % ════════════════════════════════════════════════════════════════════════
  % 底部：目录结构与依赖
  % ════════════════════════════════════════════════════════════════════════

  \draw[boundary] (0,2.9) -- (17.9,2.9);

  \node[h3,anchor=west] at (0,2.3) {项目结构};

  % 目录树（简洁版）
  \node[body,anchor=north west, font=\ttfamily\scriptsize, text=Muted] at (0,1.9) {
    \begin{minipage}[t]{4.2cm}
      \color{Ink}svg-studio/\\
      \hspace{1ex}\color{Teal}src/\\
      \hspace{2ex}main.ts\quad\color{Muted}12 场景\\
      \hspace{2ex}style.css\\
      \hspace{1ex}\color{Teal}scripts/\\
      \hspace{2ex}capture.mjs\\
      \hspace{1ex}\color{Teal}out/\quad\color{Muted}PNG 产物\\
      \hspace{1ex}catalog.json\quad\color{Muted}场景元数据\\
      \hspace{1ex}package.json\\
      \hspace{1ex}tsconfig.json
    \end{minipage}
  };

  % 依赖
  \node[h3,anchor=west] at (5.4,2.3) {核心依赖};
  \node[body,anchor=north west] at (5.4,1.9) {
    \begin{minipage}[t]{4.0cm}
      \color{Ink}\textbf{paper 0.12}\\\color{Muted}矢量几何操作\\[1.2mm]
      \color{Ink}\textbf{roughjs 4.6}\\\color{Muted}手绘质感生成\\[1.2mm]
      \color{Ink}\textbf{vite 8.x}\\\color{Muted}开发与构建\\[1.2mm]
      \color{Ink}\textbf{playwright-core}\\\color{Muted}无头浏览器测试
    \end{minipage}
  };

  % 构建命令
  \node[h3,anchor=west] at (10.0,2.3) {构建与测试};
  \node[body,anchor=north west, font=\ttfamily\scriptsize] at (10.0,1.9) {
    \begin{minipage}[t]{7.8cm}
      \color{Gold}\$ \color{Ink}npm install\\
      \color{Gold}\$ \color{Ink}npm test\\
      \color{Muted}├── typecheck (tsc --noEmit)\\
      \color{Muted}└── render (build + capture.mjs)\\
      \color{Muted}\hspace{3.5ex}├── 12 场景全量截图\\
      \color{Muted}\hspace{3.5ex}├── 交互契约验证\\
      \color{Muted}\hspace{3.5ex}└── 像素阈值断言
    \end{minipage}
  };

  % 底部脚注
  \node[note,anchor=west] at (0,0.1)
    {SVG Studio · 十二个透明画布研究 · 全部场景支持指针交互与透明背景导出};
  \node[note,anchor=east] at (17.9,0.1) {矢量即媒介 · 代码即图纸};

\end{tikzpicture}
\end{document}
